% 第15回　キャリア講習 2（記入例・模範解答）
\session{15}{3}{12月23日（水）}{キャリア講習 2}{}{}

\instr{第14回の内容を踏まえて参加します。配布資料があれば持参してください。}
\instr{\textcolor{ans}{講習の内容が未定のため、記入例は「書き方の例」として示す。}}

\afillin{講習のテーマ・講師}{（例）自分の強みを言葉にする／キャリアセンターの講師}

\work{メモ}{講習の内容}
\alines{8}{（例）自己分析のワーク：これまでの経験を「状況・課題・行動・結果」で書き出す。\par
・強みは「得意なこと」ではなく「繰り返し発揮してきた行動」から見つける\par
・ペアで経験を話し、相手から見た自分の強みを言ってもらう\par
・自己PRは、具体的な場面と数字・事実で裏づける\par
・インターンシップの探し方と、参加前に準備すること\par
・第14回の講師の話（「調べる・まとめる・説明する」）との共通点に気づいた}

\work[60mm]{整理}{学んだこと3つ}
\begin{wstabh}{colspec={Q[c,wd=6mm,m] X[2,l,m] X[3,l,m]},row{2-Z}={ht=16mm}}
  & 学んだこと & 第14回とのつながり・自分のキャリアに引きつけると \\
  1 & \af{強みは「繰り返し発揮してきた行動」から見つける} & \af{この授業では、AIの案を「選ぶ基準」を持って判断し続けた。これを強みとして説明できそう} \\
  2 & \af{自己PRは、具体的な場面と事実で裏づける} & \af{第5回の「根拠を示して反論する」と同じ。AI活用記録が、自分の行動の証拠になる} \\
  3 & \af{他者から見た強みは、自分の思い込みと違うことがある} & \af{ペアで「人の話を整理するのがうまい」と言われた。グループのリーダー経験とつながる} \\
\end{wstabh}

\work{行動}{これから半年で取り組むこと}
\alines{3}{・春休みに、関心のある業界のインターンシップに1つ参加する（2月までに応募）。\par
・この授業のレポートと提案資料を見直し、自己PRに使える場面を3つ書き出す。\par
・キャリアセンターの面談を1回予約し、自己分析の結果を話してみる。}

\begin{point}
  \item 講習の内容が決まり次第、記入例を差し替える。ここでは、第14回とのつながりと、具体的な行動計画の書き方の目安を示した。
  \item 「これから半年で取り組むこと」は、期限と行動が具体的か（いつまでに、何をするか）を見る。
  \item この授業の成果物（レポート、提案、AI活用記録）を自己PRの材料として使えることに気づかせると、授業全体の振り返りにもなる。
\end{point}
